\documentclass[11pt]{article}

%=====================================================================
%  Companion note: "The OSGS residual projection: constrained vs.
%  unconstrained finite-element spaces, and why the constrained
%  choice costs optimality."
%
%  This is a standalone pedagogical companion to the manuscript
%    "A stabilized finite element method for incompressible, inertial
%     flows in inhomogeneous porous media"
%    (Casas, Gonzalez-Usua, Codina, de-Pouplana).
%
%  It introduces NO new mathematics. Its only purpose is to explain,
%  from the ground up, why the orthogonal-subscale (OSGS) residual
%  projection is computed onto the finite-element spaces WITHOUT their
%  Dirichlet constraints, having first made the constrained choice look
%  natural and even preferable.
%
%  Cross-references of the form "(paper, eq:Label)" point to the
%  manuscript's own equation labels; they are typeset as plain text
%  and are NOT hyperlinked (this note cannot \ref labels living in
%  another document). Internal \cref/\eqref refer only to labels
%  defined inside this note.
%
%  Compiles with pdflatex / latexmk in a single pass (manual
%  thebibliography, no external .bib).
%=====================================================================

\usepackage[utf8]{inputenc}
\usepackage[margin=1in]{geometry}
\usepackage{amsmath,amssymb,amsthm,mathtools}
\usepackage{bm}
\usepackage{booktabs}
\usepackage{xcolor}
\usepackage[colorlinks=true,allcolors=blue]{hyperref}
\usepackage{cleveref}

%---------------------------------------------------------------------
%  Macros mirroring the manuscript
%---------------------------------------------------------------------
% dimensionless numbers
\newcommand{\Reyn}{\mathrm{Re}}
\newcommand{\Damk}{\mathrm{Da}}
\newcommand{\tauNS}{\tau_{1,\mathrm{NS}}}
% deviatoric-symmetric viscous projector (as in the manuscript)
\newcommand{\ViscProj}{\mathsf{P}_{\mathrm{DS}}}
% bold field symbols
\newcommand{\bu}{\boldsymbol{u}}
\newcommand{\bv}{\boldsymbol{v}}
\newcommand{\bw}{\boldsymbol{w}}
\newcommand{\bef}{\boldsymbol{f}}
\newcommand{\bx}{\boldsymbol{x}}
\newcommand{\bq}{\boldsymbol{q}}
\newcommand{\bsig}{\boldsymbol{\sigma}}
\newcommand{\ba}{\boldsymbol{a}}
\newcommand{\bz}{\boldsymbol{z}}
\newcommand{\bpi}{\boldsymbol{\pi}}
\newcommand{\bzero}{\boldsymbol{0}}
% differential operators and helpers
\newcommand{\grad}{\nabla}
\newcommand{\divg}{\nabla\!\cdot\!}
\newcommand{\dd}{\,\mathrm{d}}
\newcommand{\norm}[1]{\left\lVert #1 \right\rVert}
\DeclarePairedDelimiter{\dualp}{\langle}{\rangle}
% differential / residual operators
\newcommand{\Lop}{\mathcal{L}}
\newcommand{\Rop}{\mathcal{R}}
\newcommand{\Dop}{\mathcal{D}}
% function spaces
\newcommand{\Xspace}{\mathcal{X}}
\newcommand{\Xzero}{\mathcal{X}_{0}}
\newcommand{\Xh}{\mathcal{X}_{h}}
\newcommand{\Xhz}{\mathcal{X}_{h0}}
\newcommand{\Vspace}{\mathcal{V}}
\newcommand{\Vzero}{\mathcal{V}_{0}}
\newcommand{\Vh}{\mathcal{V}_{h}}
\newcommand{\Vhz}{\mathcal{V}_{h0}}
\newcommand{\Qspace}{\mathcal{Q}}
\newcommand{\Qzero}{\mathcal{Q}_{0}}
\newcommand{\Qh}{\mathcal{Q}_{h}}
\newcommand{\Qhz}{\mathcal{Q}_{h0}}
\newcommand{\Xtil}{\widetilde{\mathcal{X}}}
% projections
\newcommand{\Pone}{\Pi_{1}}
\newcommand{\Ponez}{\Pi_{1,0}}
\newcommand{\Poneperp}{\Pi_{1}^{\perp}}
\newcommand{\Ptwo}{\Pi_{2}}
\newcommand{\Ptwoperp}{\Pi_{2}^{\perp}}
\newcommand{\Ptauh}{\Pi_{\bm{\tau}h}}
\newcommand{\Ptilde}{\widetilde{\Pi}}
\newcommand{\Uh}{U_{h}}
\newcommand{\Util}{\widetilde{U}}
% norms
\newcommand{\normh}[1]{\lVert #1 \rVert_{h}}
\newcommand{\normtau}[1]{\lVert #1 \rVert_{\tau_{1}}}
\DeclarePairedDelimiter\vvvert{\lvert\lvert\lvert}{\rvert\rvert\rvert}
\newcommand{\tnorm}[1]{\vvvert{#1}}

%---------------------------------------------------------------------
%  Theorem-like environments (amsthm only)
%---------------------------------------------------------------------
\theoremstyle{plain}
\newtheorem{lemma}{Lemma}
\newtheorem{proposition}{Proposition}
\theoremstyle{definition}
\newtheorem{definition}{Definition}
\theoremstyle{remark}
\newtheorem*{remark}{Remark}
\newtheorem*{intuition}{Intuition}

\title{\bfseries The OSGS residual projection:\\
constrained versus unconstrained finite-element spaces,\\
and why the constrained choice costs optimality}
\author{Notes for the manuscript}
\date{\today}

\begin{document}
\maketitle

\begin{abstract}
The orthogonal-subgrid-scale (OSGS) method requires projecting the
strong residual of the momentum--mass system onto a finite-element
space. \emph{Which} space? The unknown $\Uh$ lives in the
Dirichlet-\emph{constrained} product space $\Xhz=\Vhz\times\Qhz$, so it
is natural, and analytically convenient, to project the residual onto
that same $\Xhz$. This note first builds that natural choice up
sympathetically---following Codina's own Remark~1, projecting onto the
constrained space \emph{removes} an analytical hypothesis and lets
``all the results carry over''---and then explains, from first
principles, exactly why the manuscript nonetheless projects onto the
\emph{unconstrained} spaces $\Xh=\Vh\times\Qh$, on which
$\bpi_h$ is free to be nonzero on $\partial\Omega$. The culprit is a
one-element boundary strip on which the constrained projection is
pinned to zero and cannot represent an $O(1)$ residual, seeding a
spurious boundary layer. We teach the background as we go---the VMS
scale split, the two subscale models, and $L^2$-projection
approximation theory---and we pay careful, honest attention to a
subtlety that is easy to get wrong: precisely \emph{how} optimality is
affected. Codina \emph{asserts} (Remark~1) that the constrained projection
keeps \emph{global} convergence optimal for general degree-$k$ interpolation
(only local boundary layers appear); the manuscript's footnote makes the
stronger, rate-level claim that it ``spoils the optimal $O(h^{k+1})$
convergence.'' We reconcile the two, labelling clearly what is asserted, what
is in genuine tension for $k\ge2$, and what is empirical conjecture.
\end{abstract}

\tableofcontents

%=====================================================================
\section{The setting}
%=====================================================================

\subsection{The continuous porous problem}

We consider steady incompressible flow through an inhomogeneous porous
medium on a bounded domain $\Omega\subset\mathbb{R}^d$. Writing $\bu$
for the velocity, $p$ for the pressure, $\nu$ for the kinematic
viscosity, $\alpha=\alpha(\bx)$ for the (spatially varying) porosity,
and $\bsig=\bsig(\alpha,\bu)$ for the Darcy--Brinkman--Forchheimer
resistance tensor, the strong form of the momentum and mass balances
reads, schematically,
%
\begin{subequations}\label{eq:strong}
\begin{align}
  \alpha\,\ba\cdot\grad\bu
  \;-\;\divg\!\bigl(2\alpha\nu\,\ViscProj\grad^{\mathrm s}\bu\bigr)
  \;+\;\bsig\,\bu \;+\;\alpha\,\grad p
    &= \bef && \text{in }\Omega,
    \label{eq:strongmom}\\[2pt]
  \divg(\alpha\,\bu) &= 0 && \text{in }\Omega,
    \label{eq:strongmass}
\end{align}
\end{subequations}
%
with $\bu=\bzero$ on the Dirichlet boundary $\Gamma_{\mathrm D}$
(the analysis and the numerical experiments take
$\Gamma_{\mathrm N}=\emptyset$). Here $\ba$ is the advective velocity;
in the linearised (Oseen-type) analysis it is a prescribed field with
$\divg(\alpha\ba)=0$. This is exactly the manuscript's strong momentum
and mass equations (paper, \texttt{eq:StrongMomentumEquation}) and
(paper, \texttt{eq:StrongMassEquation}); the resistance $\bsig$ is
symmetric positive semi-definite, of Forchheimer form
$\bsig=a(\alpha)+b(\alpha)\lvert\bu\rvert$ (paper,
\texttt{eq:DBFResistanceTerm}). The two dimensionless groups that
organise the regimes are the Reynolds number $\Reyn$ and the Darcy
number $\Damk$.

It is convenient to package \cref{eq:strong} as a single differential
operator acting on the pair $U\coloneqq[\bu;p]$. Let $\Lop_{\ba}$ be the
first-plus-higher-order differential operator of the system and $F$ the
data, so that \cref{eq:strong} reads $\Lop_{\ba}U=F$. The associated
\emph{strong residual} of an approximation $U_h$ is
%
\begin{equation}\label{eq:residualdef}
  \Rop U_h \;\coloneqq\; F-\Lop_{\ba}U_h ,
\end{equation}
%
the pointwise amount by which $U_h$ fails to satisfy the strong
equations. The residual is the single object around which everything
below turns. Its \emph{first-order} part---the part that dominates the
stabilisation---is
%
\begin{equation}\label{eq:firstorder}
  \bz \;\coloneqq\; \alpha\bigl(\ba\cdot\grad\bu+\grad p\bigr),
  \qquad
  \delta \;\coloneqq\; \divg(\alpha\bu),
\end{equation}
%
the convective--pressure momentum residual and the mass residual. The
reaction $\bsig\bu$ and the viscous term are, respectively, a
zeroth-order and a second-order contribution; we return to the role of
$\bsig$ in \cref{sec:annih}.

\subsection{The Galerkin spaces: constrained versus unconstrained}

The weak form is posed on
$\Xzero\coloneqq\Vzero\times\Qzero$, where
$\Vzero\coloneqq\{\bv\in H^1(\Omega)^d:\bv|_{\Gamma_{\mathrm D}}=\bzero\}$
(this reduces to $H^1_0(\Omega)^d$ when
$\Gamma_{\mathrm N}=\emptyset$) and $\Qzero\coloneqq L^2(\Omega)$
(the zero-mean subspace in the all-Dirichlet case); see (paper,
\texttt{eq:WeakForm}, \S2). The subscript $0$ records the essential
point: \emph{the trace vanishes on $\Gamma_{\mathrm D}$}.

Discretising with continuous Lagrange elements of degree $k_u$ for
velocity and $k_p$ for pressure ($k_p\le k_u+1$), we obtain \emph{two}
families of spaces, and the whole note lives in the gap between them.

\begin{definition}[Constrained and unconstrained FE spaces]
\label{def:spaces}
Let $\Vh\subset H^1(\Omega)^d$ and $\Qh\subset H^1(\Omega)\cap
C^0(\overline\Omega)$ be the \emph{unconstrained} degree-$(k_u,k_p)$
finite-element spaces---no boundary condition is imposed. Their
Dirichlet-\emph{constrained} subspaces are
\[
  \Vhz\coloneqq\Vh\cap H^1_0(\Omega)^d,
  \qquad
  \Qhz\coloneqq\Qh\cap\Qzero,
\]
i.e.\ the functions in $\Vh,\Qh$ that vanish on $\Gamma_{\mathrm D}$.
Write the product spaces
$\Xh\coloneqq\Vh\times\Qh$ (unconstrained) and
$\Xhz\coloneqq\Vhz\times\Qhz$ (constrained).
\end{definition}

This matches the manuscript verbatim: the constrained
$\Xhz\coloneqq\Vhz\times\Qhz$ appears at (paper, \S3.2) and the
appendix's line
``$\Vhz\coloneqq\Vh\cap H^1_0(\Omega)^d$, $\Qhz\coloneqq\Qh\cap\Qzero$,
$\Xhz\coloneqq\Vhz\times\Qhz$'' (paper, App.,
\texttt{oa:} setup). It is also exactly Codina's distinction between
$V_h\subset V=H^1(\Omega)$ and its subspace $V_{h,0}$ \cite{codina2008}.

\begin{remark}[Where does the discrete solution live?]
\emph{Always in the constrained space.} The Galerkin trial/test
functions must honour $\bu_h|_{\Gamma_{\mathrm D}}=\bzero$, so the
discrete unknown is $\Uh\in\Xhz$, and the Galerkin problem tests against
$V_h\in\Xhz$. This never changes. What is in question is a
\emph{different} object---the space onto which the residual is
\emph{projected}---and the temptation of this note is precisely the
(mistaken) reflex ``the residual should be projected where the unknown
lives.''
\end{remark}

%=====================================================================
\section{Subscales and the two stabilizations}\label{sec:vms}
%=====================================================================

\subsection{The variational multiscale scale split}

Standard Galerkin fails for this problem twice over: it cannot control
convection at high $\Reyn$, and equal-order $\bu$/$p$ interpolation
violates the inf--sup condition. The Variational Multiscale (VMS)
framework \cite{codina2018} repairs both from a single idea; the
stabilised method used here generalises the stationary formulation of
\cite{codina2001} and its transient orthogonal-subscale counterpart
\cite{codina2002} to an inhomogeneous porous medium. Decompose the
continuous solution space as a direct sum of what the mesh can resolve
and what it cannot,
%
\begin{equation}\label{eq:split}
  \Xzero=\Xhz\oplus\Xtil_0,
  \qquad\text{so that}\qquad
  U=\Uh+\Util,\quad \Uh\in\Xhz,\ \Util\in\Xtil_0,
\end{equation}
%
where $\Xtil_0$ is the (infinite-dimensional) \emph{subscale} or
subgrid-scale space, the complement of the finite-element space (paper,
\S3.1, \emph{Scale Splitting}; the displayed splitting there is
unlabelled). The exact solution equals the
resolved part $\Uh$ plus a subscale $\Util$ that lives below the mesh.
Testing the full weak form against resolved test functions
$V_h\in\Xhz$ gives the \emph{resolved-scale equation} (paper,
\texttt{eq:WeakFormResolvedScales}); the subscale $\Util$ appears there,
and the game is to model it and eliminate it.

\subsection{The exact subscale equation and its closure}

Restricting the subscale equation to an element $K$ and using the VMS
modelling assumptions (M1) the subscale vanishes on element boundaries,
(M2) inter-element contributions are negligible, and (M3) the operator
inverse is well approximated by a matrix of \emph{stabilisation
parameters}, $\Lop_{\ba}^{-1}|_K\approx\bm{\tau}_K$, the subscale
equation collapses to a projection statement (paper,
\texttt{eq:weak\_form\_subscales\_as\_projection}),
%
\begin{equation}\label{eq:subprojection}
  \Ptilde\bigl[\Lop\Util\bigr]=\Ptilde\bigl[\Rop U_h\bigr],
\end{equation}
%
in which $\Ptilde$ projects onto the subscale space. Solving
elementwise (paper, \texttt{eq:simplified\_weak\_form\_subscales}),
%
\begin{equation}\label{eq:submodel}
  \boxed{\;\Util|_K
  \;\approx\;\bm{\tau}_K\,\Ptilde\bigl[\Rop U_h\bigr]\big|_K\;}
\end{equation}
%
The subscale is the ($\bm\tau$-weighted) \emph{projection of the strong
residual}. Feeding \cref{eq:submodel} back into the resolved equation
yields the stabilised finite-element method (paper,
\texttt{eq:VMSWeakFormSystem}). Everything hinges on the operator
$\Ptilde$: it decides which part of the residual is fed back as a
subscale.

\subsection{ASGS versus OSGS: two choices of \texorpdfstring{$\Ptilde$}{Pi}}

There are two canonical closures.

\paragraph{ASGS (algebraic subgrid scales).}
Take the subscale space to be the space of finite-element residuals
themselves; then $\Ptilde=\mathcal I$, the identity, and the
\emph{whole} residual is retained,
$\Util|_K=\bm\tau_K\,\Rop U_h|_K$. There is no projection to speak of.

\paragraph{OSGS (orthogonal subgrid scales).}
Take the subscale space to be the $L^2$-orthogonal complement of the
finite-element space, $\Xtil=\Xhz^{\perp}$. Then $\Ptilde$ is the
\emph{orthogonal projection off} the finite-element space,
%
\begin{equation}\label{eq:osgsproj}
  \Ptilde=\mathcal I-\Ptauh
  \qquad\Longleftrightarrow\qquad
  \Util|_K=\bm\tau_K\bigl(\Rop U_h-\bpi_h\bigr)\big|_K,\quad
  \dualp{W_h,\bpi_h}=\dualp{W_h,\Rop U_h},
\end{equation}
%
where $\Ptauh$ is the ($\bm\tau$-weighted, in practice $L^2$)
projection onto the finite-element space and
$\bpi_h=\Ptauh[\Rop U_h]$ is the projected residual. The retained
subscale is the \emph{fluctuation} $\Ptilde[\Rop U_h]=(\mathcal I-\Ptauh)
[\Rop U_h]$: only the residual's component \emph{orthogonal to the
mesh} survives. This is (paper, \texttt{eq:OSGSProjection}) with the
$\bm\tau$-inner product (paper, \texttt{eq:TauInnerProduct}), simplified
to plain $L^2$ (paper, \texttt{eq:L2InnerProduct}; this simplification and
its stabilising properties are due to \cite{codina2002}). The full discrete
system is (paper, \texttt{eq:OSGSProblem}): find $\Uh\in\Xhz$ with
%
\begin{align*}
  B_{\mathrm S}(\bu;\Uh,V_h)&=L_{\mathrm S}(\bu,V_h,\bpi_h),\\
  \Util|_K&=\bm\tau_K(\bu)\bigl(\Rop U_h-\bpi_h\bigr)|_K
      \quad\forall K,\\
  \dualp{W_h,\bpi_h}&=\dualp{W_h,\Rop U_h},\\
  U&=\Uh+\Util,
\end{align*}
%
and ASGS is recovered simply by setting $\bpi_h=\bzero$.

\begin{intuition}
The projection question exists \emph{only for OSGS}. In ASGS there is no
$\bpi_h$ at all: $\Ptilde=\mathcal I$ retains the raw residual and there
is nothing to project. The entire discussion that follows---onto
\emph{which} finite-element space is $\bpi_h$ computed---is an
OSGS-only question. We will not refer to ASGS again except to say this.
\end{intuition}

Notice what \cref{eq:osgsproj} demands of us: a projection operator
$\Ptauh$ needs a target space. The unknown $\Uh$ lives in $\Xhz$; the
test functions $W_h$ defining $\bpi_h$ range over \dots\ what? That is
the one degree of freedom left, and it is the subject of the note.

%=====================================================================
\section{The natural first choice: project onto \texorpdfstring{$\Xhz$}{Xh0}}
\label{sec:natural}
%=====================================================================

Let us make the constrained choice look as attractive as it genuinely
is, because it is the choice one reaches for first, and for good
reasons.

\paragraph{Reason 1: the projected residual then lives with the unknown.}
If $W_h$ ranges over the constrained $\Xhz$, then
$\bpi_h=\Ptauh[\Rop U_h]\in\Xhz$ lives in the very same space as $\Uh$.
The subscale $\Util$, the resolved field $\Uh$, and the projected
residual $\bpi_h$ are then all objects of one type---elements of a
Dirichlet-vanishing space. There is an appealing tidiness to a method in
which every finite-dimensional quantity honours the same boundary
condition.

\paragraph{Reason 2: the analysis is strictly simpler.}
This is not a matter of taste; it is a theorem-count. In the OSGS
stability proof one needs a special test function built from the
projection of the first-order residual onto the \emph{constrained}
velocity space, $\Ponez\bz_h$ (it must be constrained to be an
admissible, boundary-vanishing test function). If the \emph{method's}
projector is \emph{also} the constrained one, then the two computable
pieces of the residual---the fluctuation $\Poneperp\bz_h$ controlled by
the Galerkin diagonal, and $\Ponez\bz_h$ supplied by the test
function---are the two complementary halves of \emph{one and the same
orthogonal projection}. They add up, by Pythagoras, to the full norm
$\normtau{\bz_h}^2$ exactly, with constant $\beta_0=1$. No extra
hypothesis is needed to recover control of the full residual.

\paragraph{Reason 3: Codina says the results carry over.}
This is stated, precisely, in \cite[Remark~1]{codina2008}. In that
work the OSGS (there ``OSS'') methods are defined by projecting the
residual onto the unconstrained $W_h$; Remark~1 records what happens if
one uses the constrained $W_{h,0}$ instead. Verbatim:

\begin{quote}\itshape
Both methods I and II could be slightly modified by projecting onto
$W_{h,0}$ in (18) instead of projecting onto $W_h$. This would simplify
the analysis presented in the following section, since the stability
condition (35) stated there would not be needed, and all the results to
be presented carry over to this case. However, even though the global
convergence is optimal, projecting onto $W_{h,0}$ leads to spurious
numerical boundary layers, similar to those found for the pressure in
classical fractional step schemes for the transient problem (see for
example [25]). Further discussion about this point can be found in [18].
\end{quote}

Read the first two sentences on their own and the constrained choice
looks like a free lunch: it \emph{removes} a hypothesis (condition~(35),
which we meet in \cref{sec:optimality}) and \emph{all the results carry
over}. If a modelling choice buys a simpler, less-assumption-heavy
theory at no cost to convergence, why would anyone do otherwise?

The answer is in the third sentence, and it is the whole point of this
note.

%=====================================================================
\section{Why the constrained projection is abandoned: the boundary strip}
\label{sec:strip}
%=====================================================================

\subsection{A refresher on \texorpdfstring{$L^2$}{L2}-projection approximation theory}

Let $\Pi_h:L^2(\Omega)\to\Vh$ be the $L^2$ projection onto the
\emph{unconstrained} degree-$k$ space: $\Pi_h g$ is the unique
$g_h\in\Vh$ with $(g-\Pi_h g,v_h)=0$ for all $v_h\in\Vh$. Because it is
an orthogonal projection it is, by construction, the \emph{best}
$L^2$-approximation,
%
\begin{equation}\label{eq:best}
  \norm{g-\Pi_h g}_{L^2(\Omega)}
  \;=\;\inf_{v_h\in\Vh}\norm{g-v_h}_{L^2(\Omega)}.
\end{equation}
%
Choosing $v_h$ to be a nodal interpolant of a smooth $g\in H^{k+1}$ and
invoking Bramble--Hilbert gives the optimal fluctuation estimate
%
\begin{equation}\label{eq:optfluct}
  \norm{g-\Pi_h g}_{L^2(\Omega)}\le C\,h^{k+1}\,\lvert g\rvert_{H^{k+1}(\Omega)}.
\end{equation}
%
\emph{The estimate \cref{eq:optfluct} is derived on the unconstrained
space}: the competitor $v_h$ in \cref{eq:best} is allowed to match $g$
up to interpolation error \emph{everywhere}, including at boundary
nodes. This freedom is exactly what will be revoked in the constrained
case.

In the OSGS method the retained subscale is precisely this kind of
fluctuation, $\Ptilde[\Rop U_h]=\Poneperp[\Rop U_h]$ with
$\Poneperp\coloneqq\mathcal I-\Pone$, but projected in the
$\bm\tau$-weighted broken norm rather than in plain $L^2$. One caveat
keeps the bookkeeping honest. The first-order residual
$\bz=\alpha(\ba\cdot\grad\bu+\grad p)$ carries \emph{one derivative
fewer} than $\bu$ (and the pressure lives only in $H^{k}$), so a plain
$L^2$ best-approximation of $\bz$ on its own would give merely
$O(h^{k-1})$, \emph{not} the $O(h^{k+1})$ one might read off
\cref{eq:optfluct}. The missing two orders are supplied by the
$\bm\tau$-weighting: on the momentum residual $\tau_1^{1/2}\sim h$, and
the paper's \emph{weighted} best-approximation lemma (paper, App.,
\texttt{oa:lem:bestapprox}, \texttt{oa:eq:WeightedApprox}), whose
$h_K^{2r}\tau_{i,K}$ weighting supplies the extra order, is what makes
the consistency estimate close at the optimal rate. In other words the
plain-$L^2$ bound \cref{eq:optfluct} is the \emph{intuition}, and the
weighted lemma is the \emph{theorem}; what matters here---and what the
constrained projection will destroy---is the same qualitative fact in
either norm: away from the boundary the fluctuation is an
\emph{interpolation-error-sized} object.

\subsection{The pinned strip}

Now switch the projector's target to the constrained $\Vhz$. Let
$\Ponez$ be the $L^2$ projection onto $\Vhz$. Every admissible
competitor $v_h\in\Vhz$ satisfies $v_h=\bzero$ on $\Gamma_{\mathrm D}$;
the projection is \emph{pinned to zero on the boundary}. But the field
being projected does not vanish there.

\begin{proposition}[Annihilation of the local projection on the strip]
\label{prop:strip}
Let $\bz=\alpha(\ba\cdot\grad\bu+\grad p)$ be the exact first-order
momentum residual \cref{eq:firstorder}, and suppose it is smooth enough
to have a well-defined trace on $\Gamma_{\mathrm D}$ with
$\bz|_{\Gamma_{\mathrm D}}\neq\bzero$ (the generic case). Let $\Omega_h$
denote the one-element strip of elements touching $\Gamma_{\mathrm D}$,
so $\lvert\Omega_h\rvert = O(h)$ on a shape-regular mesh. Then the
constrained fluctuation satisfies, pointwise,
\[
  \Ponez^{\perp}\bz \;=\; \bz-\Ponez\bz \;=\; O(1)
  \qquad\text{on }\Omega_h,
\]
and consequently, in the plain $L^2$ norm,
\begin{equation}\label{eq:half}
  \norm{\Ponez^{\perp}\bz}_{L^2(\Omega)}
  \;\ge\;\norm{\Ponez^{\perp}\bz}_{L^2(\Omega_h)}
  \;\approx\; O(1)\cdot\lvert\Omega_h\rvert^{1/2}
  \;=\; O(h^{1/2}),
\end{equation}
in place of the $O(h^{k+1})$ of \cref{eq:optfluct}.
\end{proposition}

\begin{proof}[Heuristic]
On $\Gamma_{\mathrm D}$ every $v_h\in\Vhz$ is zero, so the pointwise
defect there is forced to equal the trace itself,
$(\bz-v_h)|_{\Gamma_{\mathrm D}}=\bz|_{\Gamma_{\mathrm D}}=O(1)$,
uniformly in $h$: best approximation \cref{eq:best} cannot remove an
error the constraint hard-wires in. A continuous $\Pi_k$ function that
is $O(1)$ on $\Gamma_{\mathrm D}$ but returns toward the interior best
approximation must carry that $O(1)$ magnitude across the one boundary
layer of elements, a band of thickness $O(h)$ along a boundary of
$(d\!-\!1)$-measure $O(1)$. Hence
$\int_{\Omega_h}\lvert\Ponez^{\perp}\bz\rvert^2 \approx O(1)^2\cdot O(h)$,
which is \cref{eq:half}. The exponent $\tfrac12$ comes entirely from the
$O(h)$ \emph{measure} of the strip, not from any approximation power;
that is why it does not improve with $k$.
\end{proof}

That $\bz|_{\Gamma_{\mathrm D}}\neq\bzero$ is not exotic: on
$\Gamma_{\mathrm D}$ the velocity is Dirichlet-fixed, yet the momentum
equation forces
$\ba\cdot\grad\bu+\grad p=\alpha^{-1}(\bef+\divg(2\alpha\nu\,\ViscProj\grad^{\mathrm s}\bu)-\bsig\bu)$,
which has no reason to vanish there. So a constrained projection of the
residual is generically annihilated on a strip where the true residual
is $O(1)$: \emph{the local projection is destroyed exactly where the
subscale is supposed to live.}

\subsection{The fractional-step analogy}

This mechanism is not new; it is the spatial twin of a well-known defect
in time-splitting methods, and Codina's Remark~1 makes the analogy
explicitly (``similar to those found for the pressure in classical
fractional step schemes''). In a fractional-step / pressure-projection
scheme the velocity--pressure decoupling silently over-imposes a
boundary condition---effectively a spurious homogeneous Neumann
condition $\partial p/\partial n=0$ that the true pressure does not
satisfy---and the mismatch concentrates in a thin wall layer of width
$O(\sqrt{\nu\Delta t})$, the \emph{artificial pressure boundary layer},
which caps pressure accuracy \cite{guermond2006,codina2001pressure}. In
both cases an approximation or decomposition step forces an incorrect
boundary behaviour and the error localises in an $O(\text{width})$ strip
along $\Gamma$ where mesh (or time-step) refinement of the interior
cannot reach it. The constrained OSGS projection manufactures precisely
the spatial version of that layer.

%=====================================================================
\section{Exactly how optimality is affected --- a careful reconciliation}
\label{sec:optimality}
%=====================================================================

Here is the pedagogical payoff, and the one place where it is easy to
overclaim. There are \emph{two} statements in the literature about the
consequence of the constrained projection, and they are \emph{not} the
same statement. We put them side by side, honestly.

\subsection{(a) What Codina 2008 asserts (Remark~1): optimal global rate, local layers}

Codina's Remark~1 is a \emph{qualitative-defect} statement, not a
rate-loss statement---and it is a \emph{remark}, asserting that ``all the
results \dots\ carry over'' to the constrained variant without giving a
separate proof for it. Re-read the key clause, verbatim:

\begin{quote}\itshape
\dots\ even though the global convergence is optimal, projecting onto
$W_{h,0}$ leads to spurious numerical boundary layers \dots
\end{quote}

Global convergence \emph{remains optimal}; the constrained projection
produces \emph{spurious boundary layers}, a local, qualitative artefact.
The rate in the relevant (mesh-dependent, $\bm\tau$-weighted energy)
norm of the resolved field is not claimed to degrade. The manuscript's
own appendix transcribes this faithfully: it states only that the
constrained choice ``would simplify the analysis (assumption
\texttt{H:projection} would not be needed) but produces spurious
numerical boundary layers'' (paper, App.), and its remark on the
analysed-versus-implemented formulation asserts no rate change (paper,
App., \texttt{oa:rem:analyzed}).

\subsection{(b) The power-counting puzzle, and its resolution}

\Cref{prop:strip} produced an $O(1)$ defect on an $O(h)$-measure strip,
i.e.\ $O(h^{1/2})$ in the plain $L^2$ norm. For $k\ge1$ we have
$h^{1/2}\gg h^{k+1}$ as $h\to0$; taken at face value this would wreck an
$O(h^{k+1})$ target. How, then, can Codina keep the \emph{global} rate
optimal? Two independent mechanisms, both real, resolve the puzzle;
we flag which parts are firm and which are conjecture.

\paragraph{(i) The layer lives in the subscale, not the resolved field.}
The $O(1)$ object is $\Ptilde[\Rop U_h]=\Ponez^{\perp}[\Rop U_h]$---the
\emph{fluctuation}, hence the subscale $\Util$ and the reconstructed
total field $U=\Uh+\Util$. Codina's ``optimal global convergence'' is a
statement about the \emph{resolved} $\Uh$; the spurious layer is a
feature of the subscale/total-field reconstruction, precisely the
analogue of the fractional-step pressure layer he cites. The resolved
velocity is not itself made to carry the $O(1)$. \emph{[Firm as an
attribution of Codina's claim; the extent of leakage into $\Uh$ is the
open point of part (c).]}

\paragraph{(ii) The working norm is $\bm\tau$-weighted, and the strip is
mesh-thin.} The plain-$L^2$ count is simply the \emph{wrong norm} for
Codina's theorem: convergence is measured in the $\bm\tau$-weighted
broken norm. Redo \cref{eq:half} with the weight. On $\Omega_h$,
\begin{equation}\label{eq:weightcount}
  \normtau{\Ponez^{\perp}[\alpha X]}^2
  \;=\;\sum_{K\subset\Omega_h}\tau_{1,K}\cdot O(1)\cdot\lvert K\rvert
  \;\approx\;\tau_1\cdot O(1)\cdot h
  \quad\Longrightarrow\quad
  \normtau{\cdot}\sim(\tau_1 h)^{1/2}.
\end{equation}
In Codina's diffusion-dominated Oseen setting ($\bsig=0$),
$\tau_1\sim h^2/(c_1\alpha\nu)$, so $(\tau_1 h)^{1/2}\sim
h^{3/2}/\nu^{1/2}$. Now, Codina asserts optimal global convergence for
\emph{general} degree-$k$ interpolation---his hypothesis~H6 takes
$\bu\in H^{k+1}$, $p\in H^{k}$ with $k\ge1$, and Remark~1 says of the
constrained variant that ``all the results \dots\ carry over''; he
nowhere confines optimality to $k=1$. For $k=1$ the optimal energy rate
is $O(h)$ and the strip contribution $h^{3/2}$ is \emph{higher order},
hence invisible in the global rate: optimality is uncontroversial. For
$k\ge2$, however, the optimal energy rate is $h^{k}\ge h^2$ while this
crude $O(1)$-on-the-strip count still gives $h^{3/2}$, which \emph{would}
then dominate---a genuine tension with Codina's general-$k$ Remark~1.
\emph{[We flag this and do not pretend to resolve it. Because Remark~1 is
an informal remark with no detailed proof for the constrained variant,
we cannot tell whether the crude $O(1)$ amplitude overestimates the true
boundary defect (so that optimality survives for $k\ge2$ after all) or
whether the ``carry-over'' claim is optimistic there. What is firm is the
$k=1$ case, the existence of the tension for $k\ge2$, and the fact that
the manuscript sidesteps the whole question by projecting onto the
unconstrained space. We do \emph{not} read Codina as having restricted
optimality to $k=1$.]}

\subsection{(c) The manuscript's stronger claim, and where it rests}

The manuscript's footnote to (paper, \texttt{eq:OSGSProblem}) says
something one notch stronger. Verbatim:

\begin{quote}\itshape
In our implementation the projection \texttt{(eq:NonlinearResidualProjection)}
is computed with $W_h$ ranging over the velocity and pressure finite
element spaces \emph{without} their Dirichlet constraints, so that
$\bm\pi_h$ is not forced to vanish on $\partial\Omega$. Projecting on the
constrained space $\mathcal X_{h0}$ instead would force the boundary
degrees of freedom to vanish (the space $\mathcal X_{h0}$ carrying
homogeneous Dirichlet conditions), annihilating the local $L^2$
projection there and introducing an $O(1)$ boundary residual that spoils
the optimal $O(h^{k+1})$ convergence.
\end{quote}

The phrase ``spoils the optimal $O(h^{k+1})$ convergence'' is a
\emph{rate} claim. It goes beyond Codina's ``optimal global convergence
with boundary layers,'' and---read literally as a statement about the
global rate of the resolved field in the energy norm---it is in tension
with the paper's \emph{own} appendix, which claims only spurious layers.
This is the single sentence in the cluster that must be read carefully.

The honest reading is that the footnote's rate-level claim is \emph{not}
Codina's theorem and is \emph{not} proved in the appendix (which
explicitly declines to analyse the constrained variant). It rests
instead on the manuscript's \emph{own manufactured-solution (MMS)
evidence} in the porous/reaction setting, and it can be true
\emph{simultaneously} with Codina's theorem because the two speak about
different norms and different regimes:
%
\begin{itemize}
\item \emph{Different norm.} The MMS convergence gate measures the
  plain-$L^2$ velocity and continuity error of the \emph{resolved}
  field, via the scale-free indicators $\varepsilon_M,\varepsilon_C$.
  Plain $L^2$ is exactly the norm in which the $O(1)$-on-$O(h)$ count
  \emph{does} bite ($O(h^{1/2})$)---\emph{if} the spurious stabilisation
  forcing contaminates $\Uh$. Codina's theorem uses the $\bm\tau$-weighted
  energy norm of \cref{eq:weightcount}, which damps the strip.
\item \emph{Different regime.} Codina is $\bsig=0$ Oseen; the manuscript
  is porous, with reaction $\bsig$ and varying porosity $\alpha$ absent
  from his setting, and it probes high $\Damk$. At high $\Damk$ the exact
  solution develops a Brinkman screening layer of width
  $\ell_\sigma=(\alpha\nu/\sigma)^{1/2}$; when $h$ under-resolves
  $\ell_\sigma$, the true first-order residual $\bz$ is already steep and
  large near $\Gamma_{\mathrm D}$, and the constrained projection---still
  pinned to zero---superimposes its spurious layer on a genuine,
  under-resolved physical one, so the effective ``$O(1)$'' is
  $O(\text{large})$. It is plausible that this sharpens the qualitative
  layer into a \emph{measurable} rate loss. \emph{[Conjecture: the paper
  does not analyse the constrained variant; and note two honest
  counter-currents---(1) for constant $\sigma$ the reaction residual has
  \emph{no} fluctuation at all (\cref{sec:annih}), so $\bsig$ never drives
  the layer directly; (2) high $\Damk$ makes $\tau_1\approx1/\sigma$
  small, which \emph{suppresses} the velocity fluctuation in the energy
  norm. The rate reading therefore rests on physical-layer amplification
  plus plain-$L^2$ contamination of $\Uh$, neither of which is
  established here.]}
\item \emph{Possibly $k\ge2$.} As the power count in (b)(ii) shows,
  degradation is far more available for $k\ge2$ than for the $k=1$ case
  Codina certifies.
\end{itemize}

The disciplined conclusion: keep the footnote's rate phrasing as an
\emph{empirical, regime-specific} claim about the plain-$L^2$ MMS gate,
and do not present it as Codina's theorem or as a proven global-rate
result. Codina's qualitative statement and the paper's empirical
extension are compatible, not contradictory.

\subsection{(d) The price of the honest choice: condition (35)}

If the constrained projection is so simple analytically, what does the
\emph{unconstrained} projection---the one actually implemented and
analysed---cost? Exactly one hypothesis. In the stability proof the
Galerkin diagonal controls the fluctuation $\Poneperp\bz_h$ (now off the
\emph{unconstrained} space), while the special test function still
supplies $\Ponez\bz_h$ (necessarily off the \emph{constrained} space,
to be admissible). These two are projections onto \emph{different}
spaces, so they no longer sum to $\normtau{\bz_h}^2$ by Pythagoras. The
missing glue is the manuscript's assumption \texttt{H:projection}
(paper, \texttt{eq:ProjectionStability}), verbatim:

\begin{quote}\itshape
\emph{Projection stability.} Let $\Pi_1$ and $\Pi_{1,0}$ be the
$\tau_1$-weighted $L^2$ projections of a velocity field onto the
unconstrained and the constrained velocity spaces, and
$\Pi_1^{\perp}\coloneqq I-\Pi_1$. There is $\beta_0>0$, independent of
$h$, such that for all $V_h=[\bm v_h;q_h]\in\mathcal X_h$, writing
$\bm z_h\coloneqq\alpha(\bm a\cdot\nabla\bm v_h+\nabla q_h)$,
\end{quote}
\begin{equation}\label{eq:projstab}
  \bigl\lVert\tau_1^{1/2}\bm z_h\bigr\rVert_h
  \;\le\;\frac{1}{\beta_0}\Bigl(
    \bigl\lVert\tau_1^{1/2}\Pi_{1,0}\bm z_h\bigr\rVert_h
    +\bigl\lVert\tau_1^{1/2}\Pi_1^{\perp}\bm z_h\bigr\rVert_h\Bigr).
\end{equation}
This is exactly the Codina--Blasco-type inf--sup condition (35) of
\cite{codina2008} (the equal-order-interpolation device it descends from
originates in \cite{codinablasco1997}), which in his own notation reads
\[
  \norm{z_h}_{\tau_1}\le\frac{1}{\beta_0}
  \bigl(\norm{\Pi_{\tau_1,0}(z_h)}+\norm{\Pi_{\tau_1}^{\perp}(z_h)}_\tau\bigr),
  \qquad z_h\equiv\ba\cdot\grad\bv_h+\grad q_h .
\]
It says precisely: a bound on the constrained-projection part
$\Ponez\bz_h$ \emph{and} on the unconstrained fluctuation
$\Poneperp\bz_h$ suffices to bound the \emph{whole} residual
$\normtau{\bz_h}$---i.e.\ the component of $\bz_h$ that lies in $\Vh$ but
is $\tau_1$-orthogonal to $\Vhz$ is not independent of the other two.
(For the divergence part it holds trivially with $\beta_0=1$, because
$\bv_h=\bzero$ on $\partial\Omega$ forces $\divg\bv_h$ to have zero
mean.) It is needed precisely because $B_{\mathrm S}$ is \emph{not}
coercive; \cref{eq:projstab} is what recovers inf--sup stability from the
two controllable projections.

\begin{center}
\emph{Constrained projection:} easy analysis, no \cref{eq:projstab},
but boundary layers. \\
\emph{Unconstrained projection:} needs \cref{eq:projstab}, but clean
boundaries --- and it is the honest method that is implemented.
\end{center}

%=====================================================================
\section{Settling the question: a degree-dependent theorem}
\label{sec:settle}
%=====================================================================

We can now do more than flag the tension of \cref{sec:optimality}: in the
OSGS energy norm we can \emph{prove} the answer, and it turns out to depend
on the polynomial degree. Write $X(V)\coloneqq\ba\cdot\grad\bv+\grad q$ for
the first-order operator, so that $\bz=\alpha X(U)$ is the momentum
residual of \cref{eq:firstorder}, and let
\begin{equation}\label{eq:st:defect}
  \eta_\Pi \;\coloneqq\; \normtau{\Pi^{\perp}[\alpha X(U)]}
\end{equation}
be the \emph{consistency defect} of the method whose residual projector is
$\Pi$ (either $\Pone$, the unconstrained projection, or $\Ponez$, the
constrained one). Everything below is set in the fixed-data,
viscosity-dominated asymptotic regime that governs the convergence
\emph{rate} (\cref{sec:optimality}): $\tau_1=\Theta(h^2)$ (so
$\tau_1^{1/2}=\Theta(h)$) and $\tau_2=\Theta(1)$ (from
$\tau_2\to\nu/\alpha$; hence, by the design bound $\varepsilon\tau_2\le C_2$,
$\varepsilon=O(1)$---a fact we will need). Throughout, $\tnorm{\cdot}_{\mathrm{O}}$
is the projection-free triple norm (paper, \texttt{oa:eq:TripleNorm}),
\emph{the same} for both variants; only the method's $S_1$ changes.

\subsection{The question reduces to the size of one number}

\begin{lemma}[Reduction to the consistency defect]\label{lem:st:reduce}
Let $B_{\mathrm{OSGS}}$ be OSGS with projector $\Pi$, inf-sup stable in
$\tnorm{\cdot}_{\mathrm{O}}$ with an $h$-independent constant (for the
constrained variant this is Codina's Remark~1). Then
\begin{equation}\label{eq:st:reduce}
  \tnorm{U-U_h}_{\mathrm{O}}\;\le\;C\bigl(\,\Psi_{\mathrm{O}}(h)+\eta_\Pi\,\bigr)
  \;=\;C\bigl(\,\Theta(h^{k})+\eta_\Pi\,\bigr),
\end{equation}
where $\Psi_{\mathrm{O}}(h)=\Theta(h^{k})$ is the interpolation error in the
energy norm ($k\coloneqq k_u$).
\end{lemma}

\begin{proof}
The analyzed OSGS method is not strongly consistent; by the appendix's
consistency identity (paper, \texttt{oa:lem:consistency}), for all
$W_h\in\Xhz$,
\[
  B_{\mathrm{OSGS}}(\ba;\,U-U_h,\,W_h)=S_1(U,W_h)+S_2(U,W_h),
  \qquad S_1(U,W)=\bigl(\Pi^{\perp}[\alpha X(U)],\,\alpha X(W)\bigr)_{\tau_1}.
\]
Let $\widehat U_h\in\Xhz$ be the nodal interpolant of $U$ and
$\boldsymbol e_h\coloneqq U_h-\widehat U_h$. Inf-sup stability gives a
$W_h$ with $\tnorm{W_h}_{\mathrm{O}}=1$ and
\[
  \beta\,\tnorm{\boldsymbol e_h}_{\mathrm{O}}
  \le B_{\mathrm{OSGS}}(\boldsymbol e_h,W_h)
  = \underbrace{B_{\mathrm{OSGS}}(U-\widehat U_h,W_h)}_{(\mathrm I)}
  \;-\;\underbrace{\bigl(S_1(U,W_h)+S_2(U,W_h)\bigr)}_{(\mathrm{II})},
\]
using $B_{\mathrm{OSGS}}(U_h,W_h)=L(W_h)$ and the identity above. We do
\emph{not} invoke a general continuity of $B_{\mathrm{OSGS}}$ (it fails with
an $h$-independent constant, the convective term
$(\bw,\alpha\ba\cdot\grad\bu)$ carrying an $h^{-1}$). Instead, term (I) is
bounded by the appendix's interpolation lemma (paper,
\texttt{oa:lem:interpolation}), valid precisely because
$U-\widehat U_h$ is a \emph{small, boundary-vanishing} interpolation
error: $|(\mathrm I)|\le C\,\Psi_{\mathrm{O}}(h)\,\tnorm{W_h}_{\mathrm{O}}$.
Term (II): by Cauchy--Schwarz in $(\cdot,\cdot)_{\tau_1}$ and the full-residual
term of the triple norm, $\normtau{\alpha X(W_h)}\le\tnorm{W_h}_{\mathrm{O}}$,
so $|S_1(U,W_h)|\le\eta_\Pi\,\tnorm{W_h}_{\mathrm{O}}$; the mass part $S_2$ is
$O(h^{k_p+1})$ (the pressure carries no Dirichlet data, $\Qhz$ being only
zero-mean, so it is unaffected by the constraint). Collecting and adding
$\tnorm{U-\widehat U_h}_{\mathrm{O}}=\Theta(h^k)$ gives \cref{eq:st:reduce}.
\end{proof}

So the whole question is: \emph{how big is $\eta_\Pi$?}

\subsection{A sharp two-sided bound on the constrained defect}

\begin{lemma}[The constrained defect is $\Theta(h^{3/2})$, the unconstrained one $O(h^{k+1})$]
\label{lem:st:defect}
Assume $\bz|_{\Gamma_{\mathrm D}}=\alpha X(U)|_{\Gamma_{\mathrm D}}\neq\bzero$
on a subset of $\Gamma_{\mathrm D}$ of positive $(d\!-\!1)$-measure (the
generic case). Then, on a shape-regular quasi-uniform family in the
viscous regime,
\begin{equation}\label{eq:st:sharp}
  \eta_0\coloneqq\normtau{\Ponez^{\perp}[\alpha X(U)]}=\Theta\bigl(h^{3/2}\bigr),
  \qquad\text{while}\qquad
  \eta_1\coloneqq\normtau{\Poneperp[\alpha X(U)]}=O\bigl(h^{k+1}\bigr).
\end{equation}
The exponent $3/2$ is \emph{independent of $k$}; only the hidden constant
$c_k>0$ decreases (mildly) with $k$.
\end{lemma}

\begin{proof}
\emph{Upper bound.} By the minimizing property of the $\tau_1$-orthogonal
projection (paper, \texttt{oa:eq:BestApprox}), $\eta_0\le\normtau{\alpha X(U)-w_h}$
for every $w_h\in\Vhz$. Take $w_h$ the \emph{constrained} Scott--Zhang
interpolant, which reproduces the homogeneous nodal boundary values. Split
the mesh into the interior and the one-element boundary strip
$\Omega_h\coloneqq\bigcup\{K:\overline K\cap\Gamma_{\mathrm D}\neq\emptyset\}$,
$|\Omega_h|=O(h)$. On interior elements the constraint is inactive and the
standard estimate gives $\sum_{K}\tau_{1,K}\alpha_K^2h_K^{2k}\|\cdot\|_{H^k(K)}^2
\lesssim (h^2/\nu)\,h^{2k}=h^{2k+2}/\nu$. On a strip element $K$ the competitor
must vanish on the face $F=\overline K\cap\Gamma_{\mathrm D}$ while
$\alpha X(U)|_F\neq\bzero$; no polynomial vanishing on $F$ can match a nonzero
trace, so $\|\alpha X(U)-w_h\|_{L^2(K)}^2\ge c_k\,|K|$ with $c_k>0$ (a
subspace-obstruction / Markov fact: a fixed-degree polynomial cannot
form a boundary layer thinner than $O(h)$). Summing over the
$O(h^{-(d-1)})$ strip elements, each weighted by $\tau_{1,K}=\Theta(h^2/\nu)$,
\[
  \sum_{K\subset\Omega_h}\tau_{1,K}\|\alpha X(U)-w_h\|_{L^2(K)}^2
  \;\lesssim\;\frac{h^2}{\nu}\cdot O(1)\cdot\underbrace{h^{d}\cdot h^{-(d-1)}}_{|\Omega_h|=O(h)}
  \;=\;\frac{h^3}{\nu}.
\]
Since $2k+2>3$ for every $k\ge1$, the strip dominates and
$\eta_0\le C\,h^{3/2}\nu^{-1/2}$, with \emph{no} $k$ in the exponent. The same
computation with the \emph{un}constrained Scott--Zhang interpolant (which
matches the trace) has no strip penalty and reproduces the appendix's weighted
best approximation (paper, \texttt{oa:lem:bestapprox}), giving
$\eta_1=O(h^{k+1})$.

\emph{Lower bound (the strip penalty is real).} It suffices to exhibit one
configuration. On a uniform 1D mesh of $(0,1)$, size $h$, with linear elements,
$\tau_1$ is elementwise \emph{constant} (so the $\tau_1$-weighted
projection is the plain $L^2$ projection) \emph{and} of magnitude
$\tau_1=\Theta(h^2)$. Project $g\equiv1$ (smooth, $g|_{\partial\Omega}=1\neq0$)
onto $\Vhz$. With the interior hat basis the normal equations read
$\tfrac16 c_{i-1}+\tfrac23 c_i+\tfrac16 c_{i+1}=1$; the constant $c_i\equiv1$
solves every interior row \emph{exactly}, so the nodal error $e_i\coloneqq1-c_i$
satisfies the homogeneous recurrence $e_{i-1}+4e_i+e_{i+1}=0$ with
$e_0=e_n=1$. Its characteristic roots are $-2\pm\sqrt3$; the decaying mode has
ratio $2-\sqrt3\approx0.27$, so the error is $\Theta(1)$ on $O(1)$ elements
(width $O(h)$) and geometrically negligible inside. On the first element
$\int_0^h(g-\Ponez g)^2=\tfrac h3\bigl(1-(2-\sqrt3)+(2-\sqrt3)^2\bigr)=\Theta(h)$,
whence $\|g-\Ponez g\|_{L^2}=\Theta(h^{1/2})$ and, since this projection is the
exact minimizer, nothing is smaller. Therefore
$\eta_0=\tau_1^{1/2}\|g-\Ponez g\|_{L^2}=\Theta(h)\cdot\Theta(h^{1/2})=\Theta(h^{3/2})$.
The same one-element layer along a flat boundary patch, with
$O(h^{-(d-1)})$ faces and $|\Omega_h|=O(h)$, reproduces $\Theta(h^{3/2})$ in
$d$ dimensions.
\end{proof}

Raising $k$ helps only in the interior; it \emph{cannot} repair the trace
obstruction on the boundary, where a function pinned to $\bzero$ is asked to
represent an $O(1)$ field. That single asymmetry between $\Vh$ and $\Vhz$ is
the entire effect.

\subsection{The degree-dependent verdict}

\begin{proposition}[$k=1$: optimality is preserved]\label{prop:st:preserve}
For equal-order linear elements ($k=1$) the constrained-projection
OSGS retains the optimal energy rate,
$\tnorm{U-U_h}_{\mathrm{O}}=\Theta(h)$.
\end{proposition}

\begin{proof}
\Cref{lem:st:reduce,lem:st:defect} give
$\tnorm{U-U_h}_{\mathrm{O}}\le C(h+\eta_0)=C\bigl(h+\Theta(h^{3/2})\bigr)$. As
$3/2>1$, $\eta_0=\Theta(h^{3/2})=o(h)$ sits \emph{strictly beneath} the
interpolation floor $\Theta(h)$ and is asymptotically invisible. Hence
$\tnorm{U-U_h}_{\mathrm{O}}=\Theta(h)$.
\end{proof}

The negative half needs one natural hypothesis. In the linearized (Oseen /
Newton) problem the advecting field $\ba$ is a velocity iterate, so it is
legitimate---though not automatic in the appendix---to assume
\begin{equation*}
  \tag{$\mathrm{H}_\Gamma$}
  \ba=\bzero\ \text{on}\ \Gamma_{\mathrm D}\quad(\text{the advecting field obeys no-slip}).
\end{equation*}
Under ($\mathrm{H}_\Gamma$), on $\Gamma_{\mathrm D}$ the boundary residual is a
\emph{gradient}, $\bz|_{\Gamma_{\mathrm D}}=\alpha\grad p|_{\Gamma_{\mathrm D}}$
(the convective part vanishes), and on the strip interior $\alpha\ba\cdot\grad\bu=O(h)$
is subdominant. This is what lets a curl-free test field align with $\bz$.

\begin{proposition}[$k\ge2$: optimality is spoiled, with a matching lower bound]
\label{prop:st:spoil}
Assume ($\mathrm{H}_\Gamma$), the genericity of \cref{lem:st:defect}, and the
viscous regime ($\tau_1=\Theta(h^2)$, $\tau_2=\Theta(1)$, hence
$\varepsilon=O(1)$). For $k\ge2$,
\begin{equation}\label{eq:st:spoil}
  \tnorm{U-U_h}_{\mathrm{O}}=\Theta\bigl(h^{3/2}\bigr)\;<\;\Theta(h^{k}),
\end{equation}
so the constrained projection genuinely degrades the energy rate to $3/2$.
\end{proposition}

\begin{proof}
The upper bound is \cref{lem:st:reduce,lem:st:defect}:
$\tnorm{U-U_h}_{\mathrm{O}}\le C(h^k+\eta_0)=O(h^{3/2})$ since $3/2<k$. For the
matching lower bound, use a \emph{pressure} test function
$V^{\ast}=[\bzero;q^{\ast}]$ with $q^{\ast}\in\Qhz$ a strip-localized
finite-element bump chosen so that $\alpha\grad q^{\ast}=\bz$ on
$\Omega_h$ up to lower order---admissible because, by ($\mathrm{H}_\Gamma$),
$\bz$ is a gradient there, and because the pressure is \emph{not}
Dirichlet-constrained (so $q^{\ast}$ is free on $\Gamma_{\mathrm D}$).

\emph{Norm of the vehicle.} With $\bv^{\ast}=\bzero$ the viscous, reaction and
mass terms of $\tnorm{V^{\ast}}_{\mathrm{O}}$ vanish, leaving
$\varepsilon\|q^{\ast}\|^2=O(1)\cdot h^3=O(h^3)$ (here $\varepsilon=O(1)$ is
essential; a gradient-$O(1)$ bump of width $h$ has amplitude $O(h)$, so
$\|q^{\ast}\|^2\sim h^3$) and
$\normtau{\alpha X(V^{\ast})}^2=\normtau{\alpha\grad q^{\ast}}^2
\sim\tau_1\cdot O(1)\cdot|\Omega_h|=h^2\cdot h=h^3$. Thus
$\tnorm{V^{\ast}}_{\mathrm{O}}=\Theta(h^{3/2})$.

\emph{The pairing does not cancel.} By idempotent self-adjointness of
$\Ponez^{\perp}$ and $\alpha\grad q^{\ast}=\bz$ on $\Omega_h$,
\[
  S_1^{0}(U,V^{\ast})=\bigl(\Ponez^{\perp}\bz,\ \alpha\grad q^{\ast}\bigr)_{\tau_1}
  =\bigl(\Ponez^{\perp}\bz,\ \bz\bigr)_{\tau_1}+\text{(l.o.)}
  =\normtau{\Ponez^{\perp}\bz}^2+\text{(l.o.)}=\eta_0^2=\Theta(h^3),
\]
the lower-order correction controlled by $\eta_0\cdot o(h^{3/2})=o(h^3)$
because $\Ponez^{\perp}\bz$ is strip-concentrated (the exponential profile
of \cref{lem:st:defect}); and $S_2(U,V^{\ast})=0$.

\emph{Restricted continuity, then the bound.} Because $\bv^{\ast}=\bzero$, the
non-continuous convective term $(\bv,\alpha\ba\cdot\grad\bu)$ is absent,
and one checks termwise---using $\tau_2=\Theta(1)$ for the divergence pairing
$\normh[\tau_2^{-1}]{q^{\ast}}\le(\alpha/\nu)^{1/2}\|q^{\ast}\|\sim h^{3/2}$---that
$|B^{0}_{\mathrm{OSGS}}(W,V^{\ast})|\le C\tnorm{W}_{\mathrm{O}}\tnorm{V^{\ast}}_{\mathrm{O}}$
for \emph{this} pair. Expanding as in \cref{lem:st:reduce} with
$\boldsymbol e_h=U_h-\widehat U_h$,
\[
  B^{0}_{\mathrm{OSGS}}(\boldsymbol e_h,V^{\ast})
  =\underbrace{B^{0}_{\mathrm{OSGS}}(U-\widehat U_h,V^{\ast})}_{\le\,C\,h^{k}\cdot h^{3/2}=h^{k+3/2}}
  \;-\;\underbrace{S_1^{0}(U,V^{\ast})}_{\Theta(h^3)}
  \;=\;-\Theta(h^3)\quad(k\ge2,\ \text{since }h^{k+3/2}=o(h^3)).
\]
Restricted continuity then gives
$\tnorm{\boldsymbol e_h}_{\mathrm{O}}\ge |B^{0}_{\mathrm{OSGS}}(\boldsymbol e_h,V^{\ast})|/(C\tnorm{V^{\ast}}_{\mathrm{O}})
=\Theta(h^3)/\Theta(h^{3/2})=\Theta(h^{3/2})$, and
$\tnorm{U-U_h}_{\mathrm{O}}\ge\tnorm{\boldsymbol e_h}_{\mathrm{O}}-\tnorm{U-\widehat U_h}_{\mathrm{O}}
\ge c\,h^{3/2}-C\,h^{k}=\Theta(h^{3/2})$ for $k\ge2$. With the upper bound this
is two-sided.
\end{proof}

\begin{intuition}
The boundary blemish has a \emph{fixed} size $\Theta(h^{3/2})$, set by the
geometry of the $O(h)$ strip, not by the elements; the target it must beat is
the shrinking $h^{k}$. For $k=1$ the target $h$ is coarser than $h^{3/2}$, so
the blemish hides beneath it; for $k\ge2$ the target $h^{k}$ is finer than
$h^{3/2}$, so the blemish sticks out and \emph{becomes} the rate. Optimality
survives iff $3/2>k$, i.e.\ iff $k=1$.
\end{intuition}

\subsection{The tension, resolved}

\Cref{prop:st:preserve,prop:st:spoil} settle \cref{sec:optimality} in the
energy norm: Codina's Remark~1---``all the results carry over''---is
\emph{exactly right for $k=1$ and false for $k\ge2$}. There is no
contradiction with anything Codina proves: his consistency estimate
(\cite[Lemma~3]{codina2008}) is derived for the \emph{unconstrained} method
via a best-approximation competitor that \emph{matches the boundary
trace}---the very step that fails on $\Vhz$---and his optimal $Q_1$/$Q_2$
numerics are that unconstrained method; the constrained variant is supported
only by the informal Remark~1 and is never analysed or tested for $k\ge2$.
Symmetrically, the manuscript's footnote is right for the higher-order
elements ($\mathbb P_2$) it actually uses.

Two honest caveats bound the claim. First, the regime is load-bearing:
the argument uses $\tau_2=\Theta(1)$ (hence $\varepsilon=O(1)$); were
$\tau_2$ to scale like $h^2$ the pressure vehicle would fail, so the verdict
is tied to the viscous-dominated asymptotics that fix the rate. Second,
the theorem is in the \emph{energy} norm. The manufactured-solution gate
and the footnote's ``$O(h^{k+1})$'' measure the plain-$L^2$ velocity
error, and there the story is subtler: the defect enters the duality
(Aubin--Nitsche) estimate \emph{quadratically}
($\eta_0\cdot\eta_0^{\mathrm{dual}}=\Theta(h^3)$), so clean theory would leave
the $k=1$ $L^2$ rate optimal and only partly degrade $k\ge2$. We do not bank
on that recovery---the duality for an inconsistent, non-coercive form is
exactly the open route of the appendix, and this codebase's own MMS evidence
reports that the constrained projection breaks the $O(h^{k+1})$ gate in
practice. The energy-norm statement above is the part that is
\emph{proved}; the plain-$L^2$ verdict for $k\ge2$ is decided cleanly by
a single numerical experiment---project onto $\Vhz$ and measure the
$\mathbb P_2$ rate---which we recommend as the definitive check.

%=====================================================================
\section{A caution: two different ``annihilations''}\label{sec:annih}
%=====================================================================

The word \emph{annihilation} appears twice in this cluster with
completely unrelated meanings, and conflating them would be a serious
error.

\paragraph{The GOOD annihilation (intended).}
The OSGS appendix's ``Exact annihilations'' lemma (paper, App.,
\texttt{oa:lem:annihilation}) states, for constant $\sigma$ and
$\varepsilon$ and all $\Uh,V_h\in\Xhz$:
\begin{quote}\itshape
$\Pi_1^{\perp}(\sigma\bm u_h)=\bm 0$ and $\Pi_2^{\perp}(\varepsilon
p_h)=0$ (reactive and compressibility parts of the residual have no
fluctuation); \dots\ $(\Pi_1^{\perp}z,\sigma\bm v_h)_{\tau_1}=0$ and
$(\Pi_2^{\perp}w,\varepsilon q_h)_{\tau_2}=0$ for any $z,w$ \dots
\end{quote}
The mechanism: $\sigma\bu_h\in\Vh$ already, and the \emph{unconstrained}
fluctuation $\Poneperp=\mathcal I-\Pone$ annihilates \emph{anything
already in the projection space} $\Vh$. This is a desirable structural
fact---the reactive and compressibility residuals carry no subscale, so
OSGS keeps the full $\norm{\sigma^{1/2}\bv}$ norm and needs no reaction
renormalisation. \emph{Constancy of $\sigma$ is essential}: a variable
$\sigma$ would carry $\sigma\bu_h$ out of $\Vh$.

\paragraph{The BAD annihilation (this note).}
\Cref{prop:strip}'s ``annihilation of the local projection'' is the
destruction of $\Ponez\bz$ on the boundary strip because
\emph{constrained} functions must vanish there. It has nothing to do
with residual components lying in a space; it is a defect \emph{caused
by} the Dirichlet constraint on the projection target.

\begin{center}\itshape
Same word, opposite valence. The good annihilation is
$\Poneperp$ killing what is \emph{in} $\Vh$; the bad annihilation is the
\emph{constraint} killing $\Ponez$ on $\Gamma_{\mathrm D}$. Never
conflate them.
\end{center}

%=====================================================================
\section{Summary}\label{sec:summary}
%=====================================================================

\begin{center}
\renewcommand{\arraystretch}{1.25}
\begin{tabular}{@{}p{0.30\linewidth} p{0.30\linewidth} p{0.30\linewidth}@{}}
\toprule
 & \textbf{Constrained proj.\ onto $\Xhz$} & \textbf{Unconstrained proj.\ onto $\Xh$} \\
\midrule
Projected residual $\bpi_h$ lives in
  & $\Xhz$ (vanishes on $\Gamma_{\mathrm D}$)
  & $\Xh$ (free on $\Gamma_{\mathrm D}$) \\
Condition \cref{eq:projstab} / (35)
  & not needed ($\beta_0=1$, Pythagoras)
  & required (Codina--Blasco inf--sup) \\
Boundary behaviour of the subscale
  & $O(1)$ on the $O(h)$ strip (pinned)
  & $O(h^{k+1})$ interpolation fluctuation \\
Codina 2008 verdict (Remark~1, asserted)
  & optimal \emph{global} rate $+$ spurious \emph{local} boundary layers (general $k$; but see the $k\ge2$ tension, \S\ref{sec:optimality}(b))
  & optimal, clean boundaries \\
This paper's claimed verdict (empirical/conjectural)
  & rate-level $O(h^{k+1})$ \emph{claimed} spoiled in the porous/high-$\Damk$ plain-$L^2$ MMS gate --- regime-specific, not proved for the constrained variant
  & optimal $O(h^{k+1})$ MMS convergence \\
Implemented \& analysed here?
  & no
  & \textbf{yes} \\
\bottomrule
\end{tabular}
\end{center}

\medskip

\noindent The constrained projection is the natural first reach and the
easier theory---it even removes hypothesis \cref{eq:projstab}. It is
abandoned for one reason: on the one-element strip touching
$\Gamma_{\mathrm D}$ it is pinned to zero and cannot represent the $O(1)$
first-order residual, annihilating the local projection and seeding a
spurious boundary layer (the spatial twin of the fractional-step
pressure layer). Codina \emph{asserts} (Remark~1) that this keeps the
\emph{global} rate optimal for general $k$ while polluting the boundary
(with a genuine, unresolved tension at $k\ge2$, \S\ref{sec:optimality}(b));
the manuscript's stronger, rate-level
phrasing is an empirical extension resting on its own porous MMS, and is
best read as a claim about the plain-$L^2$ gate rather than as Codina's
energy-norm theorem. The price of the honest unconstrained
projection---and it is a fair price---is the projection-stability
condition \cref{eq:projstab}, which glues the constrained projection and
the unconstrained fluctuation back into control of the full residual.

%=====================================================================
\begin{thebibliography}{9}
%=====================================================================

\bibitem{codina2008}
R.~Codina,
\emph{Analysis of a stabilized finite element approximation of the Oseen
equations using orthogonal subscales},
Applied Numerical Mathematics \textbf{58}(3), 264--283, 2008.

\bibitem{codina2002}
R.~Codina,
\emph{Stabilized finite element approximation of transient
incompressible flows using orthogonal subscales},
Computer Methods in Applied Mechanics and Engineering
\textbf{191}(39--40), 4295--4321, 2002.

\bibitem{codina2001}
R.~Codina,
\emph{A stabilized finite element method for generalized stationary
incompressible flows},
Computer Methods in Applied Mechanics and Engineering
\textbf{190}(20--21), 2681--2706, 2001.

\bibitem{codinablasco1997}
R.~Codina and J.~Blasco,
\emph{A finite element formulation for the Stokes problem allowing equal
velocity--pressure interpolation},
Computer Methods in Applied Mechanics and Engineering
\textbf{143}, 373--391, 1997.

\bibitem{codina2018}
R.~Codina, S.~Badia, J.~Baiges and J.~Principe,
\emph{Variational multiscale methods in computational fluid dynamics},
in \emph{Encyclopedia of Computational Mechanics}, 2nd ed., Wiley,
pp.~1--28, 2018.

\bibitem{codina2001pressure}
R.~Codina,
\emph{Pressure stability in fractional step finite element methods for
incompressible flows},
Journal of Computational Physics \textbf{170}(1), 112--140, 2001.

\bibitem{guermond2006}
J.-L.~Guermond, P.~Minev and J.~Shen,
\emph{An overview of projection methods for incompressible flows},
Computer Methods in Applied Mechanics and Engineering
\textbf{195}(44--47), 6011--6045, 2006.

\end{thebibliography}

\end{document}